\documentclass[aspectratio=169,10pt]{beamer}
% ═══════════════════════════════════════════════════════════════════════════════
% SHARED PREAMBLE — Foundation Models Course (HIT)
% To be \input by each lecture file AFTER \documentclass
% ═══════════════════════════════════════════════════════════════════════════════

% ─────────────────────────────────────────────────────────────────────────────
% BEAMER THEME & BASIC SETUP
% ─────────────────────────────────────────────────────────────────────────────
\usetheme{Madrid}
\usecolortheme{default}

% Remove navigation symbols
\beamertemplatenavigationsymbolsempty

% ─────────────────────────────────────────────────────────────────────────────
% COLORS — HIT Branding
% ─────────────────────────────────────────────────────────────────────────────
\definecolor{hitblue}{RGB}{0,51,102}        % Primary: HIT Navy
\definecolor{hitgold}{RGB}{192,148,0}       % Accent: HIT Gold
\definecolor{hitlight}{RGB}{230,240,250}    % Light background

% Override Beamer palette
\setbeamercolor{primary}{fg=white,bg=hitblue}
\setbeamercolor{secondary}{fg=hitblue,bg=hitlight}
\setbeamercolor{structure}{fg=hitblue}
\setbeamercolor{alerted text}{fg=hitgold}

% ─────────────────────────────────────────────────────────────────────────────
% PACKAGES
% ─────────────────────────────────────────────────────────────────────────────
\usepackage{amsmath}
\usepackage{amssymb}
\usepackage{mathtools}
\usepackage{booktabs}
\usepackage{graphicx}
\usepackage{hyperref}
\usepackage{tikz}
\usepackage{pgfplots}
\usepackage{tcolorbox}
\usepackage{listings}
\usepackage{xcolor}
\usepackage{algorithm}
\usepackage{algpseudocode}

% ─────────────────────────────────────────────────────────────────────────────
% TikZ LIBRARIES
% ─────────────────────────────────────────────────────────────────────────────
\usetikzlibrary{shapes,arrows,positioning,calc,chains,scopes}
\pgfplotsset{compat=1.17}

% ─────────────────────────────────────────────────────────────────────────────
% CODE LISTINGS
% ─────────────────────────────────────────────────────────────────────────────
\lstset{
  basicstyle=\ttfamily\small,
  breaklines=true,
  commentstyle=\color{gray},
  keywordstyle=\color{hitblue},
  stringstyle=\color{hitgold},
  showstringspaces=false,
  backgroundcolor=\color{hitlight},
  frame=single,
  rulecolor=\color{hitblue}
}

% ─────────────────────────────────────────────────────────────────────────────
% TCOLORBOX STYLES
% ─────────────────────────────────────────────────────────────────────────────
\tcbset{
  hitbox/.style={
    colback=hitlight,
    colframe=hitblue,
    fonttitle=\bfseries,
    boxrule=1pt
  },
  alertbox/.style={
    colback=yellow!10,
    colframe=hitgold,
    fonttitle=\bfseries,
    boxrule=1.5pt
  }
}

% ─────────────────────────────────────────────────────────────────────────────
% CUSTOM COMMANDS
% ─────────────────────────────────────────────────────────────────────────────

% Placeholder for figures
\newcommand{\placeholder}[3]{
  \begin{center}
    \fbox{\begin{minipage}[c][#2][c]{#1}
      \centering\small\color{gray}[Figure: #3]
    \end{minipage}}
  \end{center}
}

% Section divider frame
\newcommand{\sectionframe}[2]{
  \begin{frame}[plain]
    \begin{tikzpicture}[remember picture,overlay]
      \fill[hitblue] (current page.south west) rectangle (current page.north east);
      \node[anchor=center, white, font=\Large\bfseries] at (current page.center) {
        \begin{tabular}{c}
          #1 \\[0.3cm]
          {\small\color{hitgold}#2}
        \end{tabular}
      };
    \end{tikzpicture}
  \end{frame}
}

% ─────────────────────────────────────────────────────────────────────────────
% LOAD CUSTOM MACROS
% ─────────────────────────────────────────────────────────────────────────────
% ═══════════════════════════════════════════════════════════════════════════════
% SHARED MACROS — Mathematical Notation for Foundation Models Course
% ═══════════════════════════════════════════════════════════════════════════════

% ─────────────────────────────────────────────────────────────────────────────
% ACTIVATION & LAYER FUNCTIONS
% ─────────────────────────────────────────────────────────────────────────────
\newcommand{\softmax}{\mathrm{softmax}}
\newcommand{\sigmoid}{\mathrm{sigmoid}}
\newcommand{\relu}{\mathrm{ReLU}}
\newcommand{\gelu}{\mathrm{GELU}}
\newcommand{\LayerNorm}{\mathrm{LayerNorm}}
\newcommand{\FFN}{\mathrm{FFN}}
\newcommand{\MHA}{\mathrm{MHA}}
\newcommand{\Attn}{\mathrm{Attention}}

% ─────────────────────────────────────────────────────────────────────────────
% ATTENTION COMPONENTS
% ─────────────────────────────────────────────────────────────────────────────
\newcommand{\query}{Q}
\newcommand{\key}{K}
\newcommand{\val}{V}
\newcommand{\dmodel}{d_{\text{model}}}
\newcommand{\dk}{d_k}
\newcommand{\nheads}{h}

% Scaled dot-product attention formula
\newcommand{\SDPA}{
  \Attn(\query, \key, \val) = \softmax\left(\frac{\query \key^T}{\sqrt{\dk}}\right) \val
}

\newcommand{\CrossAttn}{\text{CrossAttention}}

% ─────────────────────────────────────────────────────────────────────────────
% PROBABILITY & EXPECTATION
% ─────────────────────────────────────────────────────────────────────────────
\newcommand{\E}{\mathbb{E}}
\newcommand{\Prob}{\mathbb{P}}
\newcommand{\KL}{\mathrm{KL}}
\newcommand{\ELBO}{\mathrm{ELBO}}
\newcommand{\given}{\mid}

% ─────────────────────────────────────────────────────────────────────────────
% NORMS & OPERATIONS
% ─────────────────────────────────────────────────────────────────────────────
\newcommand{\norm}[1]{\left\lVert #1 \right\rVert}
\newcommand{\abs}[1]{\left| #1 \right|}
\newcommand{\inner}[2]{\langle #1, #2 \rangle}
\newcommand{\T}{^{\top}}

% ─────────────────────────────────────────────────────────────────────────────
% VECTORS & MATRICES
% ─────────────────────────────────────────────────────────────────────────────
\newcommand{\bvec}[1]{\mathbf{#1}}
\newcommand{\mat}[1]{\mathbf{#1}}
\newcommand{\iid}{\stackrel{\text{i.i.d.}}{\sim}}
\newcommand{\defeq}{\coloneq}

% ─────────────────────────────────────────────────────────────────────────────
% LANGUAGE MODELING
% ─────────────────────────────────────────────────────────────────────────────
\newcommand{\vocab}{\mathcal{V}}
\newcommand{\context}{\text{context}}
\newcommand{\token}{t}
\newcommand{\seq}{x}
\newcommand{\ptheta}{p_\theta}
\newcommand{\pref}{p_{\text{ref}}}

% ─────────────────────────────────────────────────────────────────────────────
% RLHF & DPO
% ─────────────────────────────────────────────────────────────────────────────
\newcommand{\piref}{\pi_{\text{ref}}}
\newcommand{\pitheta}{\pi_\theta}
\newcommand{\yw}{y_w}
\newcommand{\yl}{y_l}
\newcommand{\DPOloss}{\mathcal{L}_{\mathrm{DPO}}}

% ─────────────────────────────────────────────────────────────────────────────
% RETRIEVAL & RAG
% ─────────────────────────────────────────────────────────────────────────────
\newcommand{\docs}{\mathcal{D}}
\newcommand{\qvec}{q}
\newcommand{\dvec}{d}

% ─────────────────────────────────────────────────────────────────────────────
% AGENTS & REASONING
% ─────────────────────────────────────────────────────────────────────────────
\newcommand{\obs}{o}
\newcommand{\act}{a}
\newcommand{\thought}{t}
\newcommand{\tools}{\mathcal{T}}
\newcommand{\history}{h}
\newcommand{\concat}{\oplus}

% ─────────────────────────────────────────────────────────────────────────────
% SETS & LOGIC
% ─────────────────────────────────────────────────────────────────────────────
\newcommand{\R}{\mathbb{R}}
\newcommand{\N}{\mathbb{N}}
\newcommand{\Z}{\mathbb{Z}}

\endinput


% ─────────────────────────────────────────────────────────────────────────────
% TITLE & AUTHOR (can be overridden in individual lectures)
% ─────────────────────────────────────────────────────────────────────────────
\author[D. Lederman]{Dror Lederman}
\institute{Holon Institute of Technology (HIT) \\ Data Science Department}
\date{\today}

% ─────────────────────────────────────────────────────────────────────────────
% FOOTER & HEADER
% ─────────────────────────────────────────────────────────────────────────────
\setbeamertemplate{footline}{
  \leavevmode%
  \hbox{%
    \begin{beamercolorbox}[wd=0.33\paperwidth,ht=2.25ex,dp=1ex,center]{author in head/foot}%
      \usebeamerfont{author in head/foot}\insertshortauthor
    \end{beamercolorbox}%
    \begin{beamercolorbox}[wd=0.34\paperwidth,ht=2.25ex,dp=1ex,center]{title in head/foot}%
      \usebeamerfont{title in head/foot}\insertshorttitle
    \end{beamercolorbox}%
    \begin{beamercolorbox}[wd=0.33\paperwidth,ht=2.25ex,dp=1ex,right]{frame number}%
      \usebeamerfont{frame number}\insertframenumber{} / \inserttotalframenumber\hspace*{2ex}
    \end{beamercolorbox}%
  }%
  \vskip0pt%
}

\endinput


\title{Week 7: Foundations of Agentic AI}
\subtitle{Foundation Models and Agentic AI}
\author{Lecturer Name}
\institute{Holon Institute of Technology (HIT)}
\date{Semester, 2025}

\begin{document}

\begin{frame}
  \titlepage
\end{frame}

\begin{frame}{Learning Objectives}
  After this lecture you will be able to:
  \begin{itemize}
    \item Define an \textbf{AI agent} and distinguish it from a prompted LLM.
    \item Describe the \textbf{ReAct} framework and implement the Thought-Action-Observation loop.
    \item Explain \textbf{function calling} and its JSON schema mechanism.
    \item Understand how \textbf{Toolformer} taught models to self-learn tool use.
    \item Identify the key components of an \textbf{agent loop} and their interactions.
    \item Recognize the \textbf{planning modes} available to LLM agents.
  \end{itemize}
\end{frame}

\section{What Is an AI Agent?}
\sectionframe{What Is an AI Agent?}{Beyond prompted completion}

\begin{frame}{From LLM to Agent: A Spectrum}
  \begin{center}
  \begin{tabular}{lll}
    \toprule
    \textbf{System} & \textbf{Control flow} & \textbf{Environment interaction} \\
    \midrule
    Prompted LLM & Single forward pass & None \\
    Chain / Pipeline & Hardcoded multi-step & None / limited \\
    \textbf{Agent} & LLM-directed, iterative & Yes (tools, APIs, world) \\
    \bottomrule
  \end{tabular}
  \end{center}
  \vspace{0.5em}
  \begin{tcolorbox}[hitbox, title=Definition]
    An \textbf{AI agent} is a system that uses a language model as its controller to
    \textbf{perceive} observations, \textbf{reason} about what to do, \textbf{act} using tools,
    and \textbf{iterate} until a task is complete.
  \end{tcolorbox}
  \note{Ask students: how is this different from a for-loop? The key is that the LLM decides when to stop and what to do next, not the programmer.}
\end{frame}

\begin{frame}{Key Properties of an Agent}
  \begin{columns}[T]
    \column{0.52\textwidth}
      \begin{enumerate}
        \item \textbf{Goal-directed:} Pursues a specified objective across multiple steps.
        \item \textbf{Reactive:} Responds to observations from the environment.
        \item \textbf{Proactive:} Takes initiative; doesn't wait for instructions.
        \item \textbf{Persistent:} Maintains state (history, memory) across steps.
        \item \textbf{Tool-using:} Extends capabilities beyond LLM generation.
      \end{enumerate}
    \column{0.44\textwidth}
      % tikz/agent_loop.tex
% ReAct agent loop: Thought → Action → Observation
% Usage: % tikz/agent_loop.tex
% ReAct agent loop: Thought → Action → Observation
% Usage: % tikz/agent_loop.tex
% ReAct agent loop: Thought → Action → Observation
% Usage: \input{../../tikz/agent_loop}
\begin{tikzpicture}[
  node distance=1.0cm and 1.2cm,
  main/.style={rectangle, draw=hitblue, fill=hitlightblue, rounded corners=5pt,
               minimum width=2.2cm, minimum height=0.8cm, font=\small\bfseries, align=center},
  env/.style={rectangle, draw=hitgold, fill=hitgold!15, rounded corners=5pt,
              minimum width=2.4cm, minimum height=0.8cm, font=\small, align=center},
  tool/.style={rectangle, draw=gray, fill=gray!8, rounded corners=3pt,
               minimum width=1.6cm, minimum height=0.5cm, font=\scriptsize, align=center},
  arr/.style={->, >=Stealth, very thick, hitblue},
  earr/.style={->, >=Stealth, thick, hitgold},
  lbl/.style={font=\scriptsize, gray},
]

% Core loop nodes
\node[main] (llm) {LLM};
\node[main, right=2.0cm of llm] (action) {Action};
\node[main, below=1.4cm of action] (obs) {Observation};
\node[main, left=2.0cm of obs] (thought) {Thought};

% Environment
\node[env, right=1.5cm of action] (env) {Environment /\\Tools};

% Arrows: main loop
\draw[arr] (llm.east) -- node[above, lbl] {generates} (action.west);
\draw[arr] (action.south) -- node[right, lbl] {executes} (obs.north);
\draw[arr] (obs.west) -- node[below, lbl] {returns} (thought.east);
\draw[arr] (thought.north) -- node[left, lbl] {informs} (llm.south);

% Arrow: action → environment
\draw[earr] (action.east) -- node[above, lbl] {call} (env.west);
\draw[earr] (env.south) -- ++(0,-0.4) -| node[below right, lbl, pos=0.3] {result} (obs.east);

% Task input
\node[env, above=0.7cm of llm] (task) {Task / Prompt};
\draw[earr] (task) -- (llm);

% Tool examples below environment
\node[tool, below=0.3cm of env] (t1) {Search};
\node[tool, right=0.3cm of t1] (t2) {Calculator};
\node[tool, below=0.3cm of t1] (t3) {Code exec.};
\node[tool, right=0.3cm of t3] (t4) {Browser};

\end{tikzpicture}

\begin{tikzpicture}[
  node distance=1.0cm and 1.2cm,
  main/.style={rectangle, draw=hitblue, fill=hitlightblue, rounded corners=5pt,
               minimum width=2.2cm, minimum height=0.8cm, font=\small\bfseries, align=center},
  env/.style={rectangle, draw=hitgold, fill=hitgold!15, rounded corners=5pt,
              minimum width=2.4cm, minimum height=0.8cm, font=\small, align=center},
  tool/.style={rectangle, draw=gray, fill=gray!8, rounded corners=3pt,
               minimum width=1.6cm, minimum height=0.5cm, font=\scriptsize, align=center},
  arr/.style={->, >=Stealth, very thick, hitblue},
  earr/.style={->, >=Stealth, thick, hitgold},
  lbl/.style={font=\scriptsize, gray},
]

% Core loop nodes
\node[main] (llm) {LLM};
\node[main, right=2.0cm of llm] (action) {Action};
\node[main, below=1.4cm of action] (obs) {Observation};
\node[main, left=2.0cm of obs] (thought) {Thought};

% Environment
\node[env, right=1.5cm of action] (env) {Environment /\\Tools};

% Arrows: main loop
\draw[arr] (llm.east) -- node[above, lbl] {generates} (action.west);
\draw[arr] (action.south) -- node[right, lbl] {executes} (obs.north);
\draw[arr] (obs.west) -- node[below, lbl] {returns} (thought.east);
\draw[arr] (thought.north) -- node[left, lbl] {informs} (llm.south);

% Arrow: action → environment
\draw[earr] (action.east) -- node[above, lbl] {call} (env.west);
\draw[earr] (env.south) -- ++(0,-0.4) -| node[below right, lbl, pos=0.3] {result} (obs.east);

% Task input
\node[env, above=0.7cm of llm] (task) {Task / Prompt};
\draw[earr] (task) -- (llm);

% Tool examples below environment
\node[tool, below=0.3cm of env] (t1) {Search};
\node[tool, right=0.3cm of t1] (t2) {Calculator};
\node[tool, below=0.3cm of t1] (t3) {Code exec.};
\node[tool, right=0.3cm of t3] (t4) {Browser};

\end{tikzpicture}

\begin{tikzpicture}[
  node distance=1.0cm and 1.2cm,
  main/.style={rectangle, draw=hitblue, fill=hitlightblue, rounded corners=5pt,
               minimum width=2.2cm, minimum height=0.8cm, font=\small\bfseries, align=center},
  env/.style={rectangle, draw=hitgold, fill=hitgold!15, rounded corners=5pt,
              minimum width=2.4cm, minimum height=0.8cm, font=\small, align=center},
  tool/.style={rectangle, draw=gray, fill=gray!8, rounded corners=3pt,
               minimum width=1.6cm, minimum height=0.5cm, font=\scriptsize, align=center},
  arr/.style={->, >=Stealth, very thick, hitblue},
  earr/.style={->, >=Stealth, thick, hitgold},
  lbl/.style={font=\scriptsize, gray},
]

% Core loop nodes
\node[main] (llm) {LLM};
\node[main, right=2.0cm of llm] (action) {Action};
\node[main, below=1.4cm of action] (obs) {Observation};
\node[main, left=2.0cm of obs] (thought) {Thought};

% Environment
\node[env, right=1.5cm of action] (env) {Environment /\\Tools};

% Arrows: main loop
\draw[arr] (llm.east) -- node[above, lbl] {generates} (action.west);
\draw[arr] (action.south) -- node[right, lbl] {executes} (obs.north);
\draw[arr] (obs.west) -- node[below, lbl] {returns} (thought.east);
\draw[arr] (thought.north) -- node[left, lbl] {informs} (llm.south);

% Arrow: action → environment
\draw[earr] (action.east) -- node[above, lbl] {call} (env.west);
\draw[earr] (env.south) -- ++(0,-0.4) -| node[below right, lbl, pos=0.3] {result} (obs.east);

% Task input
\node[env, above=0.7cm of llm] (task) {Task / Prompt};
\draw[earr] (task) -- (llm);

% Tool examples below environment
\node[tool, below=0.3cm of env] (t1) {Search};
\node[tool, right=0.3cm of t1] (t2) {Calculator};
\node[tool, below=0.3cm of t1] (t3) {Code exec.};
\node[tool, right=0.3cm of t3] (t4) {Browser};

\end{tikzpicture}

  \end{columns}
\end{frame}

\section{The ReAct Framework}
\sectionframe{ReAct}{Synergizing Reasoning and Acting}

\begin{frame}{ReAct: Thought-Action-Observation}
  \textbf{Yao et al.\ (2023)}~\cite{yao2023react}
  \vspace{0.4em}
  \begin{tcolorbox}[hitbox, title=ReAct Trace Format]
    \textbf{Thought:} I need to find the capital of France.\\
    \textbf{Action:} Search[capital of France]\\
    \textbf{Observation:} Paris is the capital and most populous city of France.\\[4pt]
    \textbf{Thought:} Now I know the answer.\\
    \textbf{Action:} Finish[Paris]
  \end{tcolorbox}
  \vspace{0.3em}
  \textbf{Key insight:} Interleaving \emph{reasoning traces} with \emph{action steps} allows the model to:
  \begin{itemize}
    \item Ground reasoning in real observations (not hallucinations).
    \item Update the plan based on what tools return.
    \item Produce interpretable, auditable reasoning.
  \end{itemize}
\end{frame}

\begin{frame}{ReAct Algorithm}
  \begin{algorithm}[H]
    \caption{ReAct Agent Loop}
    \begin{algorithmic}[1]
      \Require Task $\tau$, LLM $\mathcal{M}$, tool set $\tools$, max steps $T$
      \State $\history \gets [\tau]$
      \For{$t = 1$ to $T$}
        \State $\thought_t \gets \mathcal{M}(\history)$ \Comment{Generate reasoning trace}
        \State $\act_t \gets \mathcal{M}(\history \concat [\thought_t])$ \Comment{Generate action}
        \If{$\act_t = \text{Finish}[y]$}
          \Return $y$
        \EndIf
        \State $\obs_t \gets \text{Execute}(\act_t, \tools)$ \Comment{Call tool, get observation}
        \State $\history \gets \history \concat [\thought_t, \act_t, \obs_t]$
      \EndFor
      \Return \text{Failed}
    \end{algorithmic}
  \end{algorithm}
\end{frame}

\begin{frame}{ReAct vs.\ Act-Only and Reason-Only}
  \begin{columns}[T]
    \column{0.50\textwidth}
      \textbf{Ablations (Yao et al.\ 2023):}
      \begin{itemize}
        \item \textbf{Act-only}: take actions without reasoning $\Rightarrow$ gets stuck, no backtracking.
        \item \textbf{CoT-only}: reason without acting $\Rightarrow$ hallucinates facts.
        \item \textbf{ReAct}: best of both — grounds reasoning in observations.
      \end{itemize}
      \vspace{0.4em}
      \textbf{Results on HotpotQA:}\\
      ReAct +10\% over CoT-only, +15\% over Act-only.
    \column{0.46\textwidth}
      \placeholder{0.95\textwidth}{3.5cm}{ReAct vs ablations on HotpotQA, Fever}
  \end{columns}
\end{frame}

\section{Tool Use and Function Calling}
\sectionframe{Tool Use and Function Calling}{Giving agents hands}

\begin{frame}{What Are Tools?}
  \begin{tcolorbox}[hitbox, title=Tool = External Capability]
    A \textbf{tool} is any external function or API that an agent can call:
    \begin{itemize}
      \item Web search (Bing, DuckDuckGo)
      \item Code execution (Python REPL, sandbox)
      \item Database query (SQL, vector search)
      \item API calls (weather, calendar, email)
      \item File read/write
      \item Calculator, unit converter
    \end{itemize}
  \end{tcolorbox}
  \vspace{0.3em}
  Tools extend the agent's \emph{knowledge} (search), \emph{precision} (calculator), and
  \emph{agency} (email send, file write).
\end{frame}

\begin{frame}{Function Calling (OpenAI / Anthropic)}
  \textbf{Function calling} = structured tool use via JSON schema.
  \vspace{0.4em}
  \begin{columns}[T]
    \column{0.52\textwidth}
      \textbf{Tool definition:}
\begin{lstlisting}[language=Python,basicstyle=\ttfamily\tiny]
{
  "name": "get_weather",
  "description": "Get current weather for a city",
  "parameters": {
    "type": "object",
    "properties": {
      "city": {
        "type": "string",
        "description": "City name"
      }
    },
    "required": ["city"]
  }
}
\end{lstlisting}
    \column{0.44\textwidth}
      \textbf{Flow:}
      \begin{enumerate}
        \item Pass tool schema to LLM.
        \item LLM decides to call \texttt{get\_weather(\{"city": "Paris"\})}.
        \item Host executes the call.
        \item Return result to LLM as a \texttt{tool} message.
        \item LLM continues generating.
      \end{enumerate}
  \end{columns}
  \note{OpenAI introduced function calling in June 2023. It standardized how models express intent to call tools in a structured, parseable way.}
\end{frame}

\begin{frame}{Toolformer: Self-Supervised Tool Learning}
  \textbf{Schick et al.\ (2023)}~\cite{schick2023toolformer}
  \vspace{0.4em}
  \begin{tcolorbox}[hitbox, title=Key Idea]
    Teach a language model to \textbf{insert API calls} in its own output, \emph{without} human
    demonstrations, by using self-supervised filtering:
    \begin{enumerate}
      \item Propose candidate API call positions via few-shot prompting.
      \item Execute the API call; measure if it \emph{reduces perplexity} on the continuation.
      \item Keep API calls that help; discard those that don't.
      \item Fine-tune on filtered data with API calls inline.
    \end{enumerate}
  \end{tcolorbox}
  \vspace{0.3em}
  \textbf{Tools:} calculator, Q\&A, Wikipedia search, machine translation, calendar.\\
  Result: small (6.7B) model with tool use outperforms much larger models on downstream tasks.
\end{frame}

\section{Agent Planning Modes}
\sectionframe{Agent Planning}{How agents decide what to do next}

\begin{frame}{Planning Modes}
  \begin{columns}[T]
    \column{0.52\textwidth}
      \textbf{1. Single-path (ReAct):}
      \begin{itemize}
        \item Generate one thought + action at each step.
        \item Greedy; cannot backtrack.
      \end{itemize}
      \vspace{0.3em}
      \textbf{2. Multi-path (Tree of Thoughts):}
      \begin{itemize}
        \item Generate and evaluate multiple candidate actions.
        \item BFS/DFS/MCTS over thought tree.
      \end{itemize}
      \vspace{0.3em}
      \textbf{3. Plan-then-execute:}
      \begin{itemize}
        \item LLM first generates a full plan as a list.
        \item Then executes each step.
        \item Brittle to mid-plan surprises.
      \end{itemize}
    \column{0.44\textwidth}
      \textbf{4. Iterative replanning:}
      \begin{itemize}
        \item Make a plan; execute step 1; observe; replan.
        \item More robust but more expensive.
      \end{itemize}
      \vspace{0.3em}
      \textbf{5. Hierarchical planning:}
      \begin{itemize}
        \item High-level planner decomposes into sub-tasks.
        \item Sub-agents execute each sub-task.
      \end{itemize}
      \vspace{0.3em}
      We will cover Tree of Thoughts and Reflexion in \textbf{Week 8}.
  \end{columns}
\end{frame}

\section{Lab / Demo}
\begin{frame}{Lab Idea: Build a ReAct Agent from Scratch}
  \begin{tcolorbox}[hitbox, title=Lab 7 — ReAct Agent]
    \textbf{Goal:} Implement the ReAct loop manually (not via a framework).\\[4pt]
    \textbf{Tools:} Python, OpenAI or Anthropic API, DuckDuckGo search API
    \begin{enumerate}
      \item Define a system prompt with Thought/Action/Observation format.
      \item Implement the loop: call LLM → parse action → execute → append observation → repeat.
      \item Tools: \texttt{search(query)}, \texttt{calculator(expr)}, \texttt{finish(answer)}.
      \item Test on 5 HotpotQA multi-hop questions.
      \item Measure: answer accuracy, number of steps, failure modes.
    \end{enumerate}
    \textbf{Deliverable:} Working code + trace analysis for 5 examples.
  \end{tcolorbox}
\end{frame}

\section{Discussion \& Summary}
\begin{frame}{Discussion Questions}
  \begin{enumerate}
    \item ReAct interleaves reasoning and acting. What happens when the model's reasoning trace is incorrect but the action happens to succeed anyway? Is an incorrect trace harmful?
    \item Function calling relies on the LLM to correctly format JSON. What happens when the model produces malformed JSON? How would you make a robust agent?
    \item Toolformer uses perplexity reduction as a signal for whether an API call was useful. What are the limitations of this proxy? Could an API call be useful but not reduce perplexity?
    \item What is the maximum number of agent steps before a system becomes unsafe to run autonomously? How should this limit be set?
  \end{enumerate}
\end{frame}

\begin{frame}{Summary}
  \begin{itemize}
    \item An \textbf{AI agent} is an LLM-controlled system that perceives, reasons, acts, and iterates.
    \item \textbf{ReAct} interleaves Thought-Action-Observation traces for grounded, auditable reasoning.
    \item \textbf{Function calling} standardizes structured tool invocation via JSON schemas.
    \item \textbf{Toolformer} showed models can learn tool use without human demonstrations.
    \item Planning modes range from greedy single-path to hierarchical multi-agent — each with trade-offs.
    \item Tools extend agents' knowledge (search), precision (calculator), and action space (APIs).
  \end{itemize}
  \vspace{0.3em}
  \textbf{Next week:} Deep dive into reasoning and planning — CoT, ToT, Reflexion, MCTS.
\end{frame}

\begin{frame}[allowframebreaks]{Suggested Reading}
  \nocite{yao2023react, schick2023toolformer, nakano2021webgpt}
  \bibliography{../../references}
\end{frame}

\end{document}
